% Skróty należy uporządkować alfabetycznie według pierwszej kolumny.
% Przy pierwszym wystąpieniu w tekście rozwiń skrót, np.:
% programowalny sterownik logiczny (PLC).
\cleardoublepage
\chapter*{\ploren{Wykaz skrótów}{List of abbreviations}}
\phantomsection
\addcontentsline{toc}{chapter}{\ploren{Wykaz skrótów}{List of abbreviations}}

\begin{longtable}{@{}>{\bfseries}p{0.18\textwidth}p{0.76\textwidth}@{}}
  \toprule
  \ploren{Skrót}{Abbreviation} & \ploren{Znaczenie}{Meaning} \\
  \midrule
  \endfirsthead
  \toprule
  \ploren{Skrót}{Abbreviation} & \ploren{Znaczenie}{Meaning} \\
  \midrule
  \endhead
  AI   & \ploren{sztuczna inteligencja (ang. artificial intelligence)}{artificial intelligence} \\
  IoT  & \ploren{Internet Rzeczy (ang. Internet of Things)}{Internet of Things} \\
  PLC  & \ploren{programowalny sterownik logiczny (ang. programmable logic controller)}{programmable logic controller} \\
  SCADA & \ploren{system nadzoru i akwizycji danych (ang. supervisory control and data acquisition)}{supervisory control and data acquisition} \\
  \bottomrule
\end{longtable}
